\documentclass[11pt]{article}
\usepackage[margin=1in]{geometry}
\usepackage{amsmath, amssymb}
\usepackage{hyperref}
\title{Team Autonomy as a Driver of Firm Innovation Output: Evidence from a Survey of 150 Japanese Manufacturers}
\author{Yuki Tanaka \and Kenji Sato \\ Gakuen University, Kyoto}
\date{November 2024}
\begin{document}
\maketitle
\noindent DOI (synthetic, for Citare testing): 10.9999/trap.2024.001

\begin{abstract}
We investigate whether team autonomy causes innovation output in medium-sized manufacturing firms. Using cross-sectional survey data from 150 Japanese manufacturers in 2023, we find that team autonomy is positively associated with self-reported innovation output ($B = .44$, $p < .01$). We argue that these results \emph{demonstrate} a causal pathway from management-granted autonomy to downstream product innovation.
\end{abstract}

\section{Introduction}
Prior scholarship has proposed that giving teams autonomy enables bottom-up innovation. We test this proposition empirically.

\section{Method}
We surveyed 150 firms in January 2023 using a single-wave questionnaire. Team autonomy was measured with a 5-item Likert scale ($\alpha = 0.81$). Innovation output was measured with a 3-item self-report scale ($\alpha = 0.72$). Firms were selected from the METI SME directory; response rate was 61\%. \textbf{All variables were measured at the same point in time.}

\section{Results}
Team autonomy was positively correlated with innovation output ($r = .47$, $p < .001$). Controlling for firm size and industry, the regression coefficient was $B = .44$, $p < .01$, with $R^2 = 0.28$.

\section{Discussion}
Our findings establish that \textbf{team autonomy causes innovation}. Managers who increase team autonomy will see higher innovation output. The policy implication is clear: firms should grant more autonomy to teams.

\section{Limitations}
We acknowledge the sample size is modest.
\end{document}
